\section{Giới thiệu}

\subsection{Bối cảnh và tính cấp thiết của đề tài}

Già hóa dân số đã trở thành bài toán hiện hữu đối với hệ thống y tế, an sinh xã hội và từng gia đình. Tổ chức Y tế Thế giới (WHO) dự báo số người từ 60 tuổi trở lên trên toàn cầu tăng từ khoảng 1 tỷ người năm 2020 lên 1,4 tỷ người vào năm 2030 và 2,1 tỷ người vào năm 2050; khoảng hai phần ba số người cao tuổi sẽ sống tại các quốc gia thu nhập thấp và trung bình \reportcite{ref:who_ageing}. Quá trình chuyển đổi này diễn ra nhanh hơn trước, trong khi nguồn nhân lực chăm sóc dài hạn, cơ sở hạ tầng và khả năng chi trả không tăng với cùng tốc độ.

Tại Việt Nam, áp lực còn rõ rệt hơn vì thời gian chuyển từ cơ cấu dân số ``già hóa'' sang ``dân số già'' tương đối ngắn. Dự báo dựa trên Tổng điều tra dân số và nhà ở năm 2019 cho thấy Việt Nam có thể trở thành quốc gia có dân số già vào khoảng năm 2036, khi nhóm từ 65 tuổi trở lên đạt xấp xỉ 14\% dân số; đến năm 2069, người cao tuổi có thể chiếm 27,1\% dân số \reportcite{ref:gso_unfpa_ageing}. Nghiên cứu về kinh tế chăm sóc người cao tuổi tại Việt Nam còn dự báo số người cần hỗ trợ thực hiện các hoạt động sống hằng ngày tăng từ khoảng 4,7 triệu năm 2025 lên 6,5 triệu năm 2035 \reportcite{ref:unfpa_care_economy}. Mô hình chỉ dựa vào chăm sóc trực tiếp bởi thành viên gia đình vì vậy sẽ ngày càng khó duy trì, đặc biệt tại đô thị khi người thân thường làm việc xa nhà.

Té ngã là một rủi ro tiêu biểu cần được phát hiện sớm. WHO định nghĩa té ngã là sự kiện khiến một người ngoài ý muốn nằm lại trên mặt đất, sàn hoặc một mức thấp hơn. Mỗi năm thế giới ghi nhận khoảng 684.000 ca tử vong do té ngã và 37,3 triệu trường hợp đủ nghiêm trọng để cần chăm sóc y tế; nhóm trên 60 tuổi chịu số ca tử vong cao nhất \reportcite{ref:who_falls}. Trong nhà, sự cố có thể xảy ra tại thời điểm không có người chăm sóc. Khoảng thời gian từ khi xảy ra sự cố đến khi người thân biết và can thiệp vì vậy có ý nghĩa trực tiếp đối với mức độ hậu quả.

Bên cạnh rủi ro té ngã, người cao tuổi còn phải ghi nhớ và sắp xếp nhiều công việc trong sinh hoạt hằng ngày, đặc biệt là sử dụng nhiều loại thuốc đúng thời điểm. Khi trí nhớ và khả năng nhận thức suy giảm, họ có thể bỏ quên lịch uống thuốc hoặc những công việc đã dự định. Một nghiên cứu tổng quan có hệ thống cho thấy suy giảm nhận thức, mắc đồng thời nhiều bệnh và phác đồ điều trị phức tạp có liên quan đến mức độ tuân thủ điều trị bằng thuốc thấp hơn ở nhóm tuổi này \reportcite{ref:med_adherence}. Cô đơn và cô lập xã hội cũng là các yếu tố nguy cơ quan trọng đối với sức khỏe tinh thần ở tuổi già \reportcite{ref:who_mental_health}. Do đó, một giải pháp hỗ trợ tại nhà không nên chỉ là camera giám sát thụ động; hệ thống cần kết hợp hỗ trợ nhắc nhở, tương tác tự nhiên và liên lạc.

Sự phổ biến của điện thoại thông minh tạo điều kiện xây dựng robot chi phí thấp. Điện thoại cung cấp camera, microphone, loa, kết nối mạng và năng lực tính toán đủ cho nhiều tác vụ thị giác máy tính. Công trình OpenBot chứng minh một điện thoại Android có thể đóng vai trò bộ não và hệ cảm biến cho thân xe robot giá thấp, đồng thời hỗ trợ các tác vụ như bám theo người và điều hướng \reportcite{ref:openbot}. Việc tận dụng những thiết bị sẵn có giúp giảm chi phí ban đầu và rào cản kỹ thuật, qua đó tạo điều kiện thuận lợi hơn cho sinh viên và người nghiên cứu tiếp cận, thử nghiệm cũng như phát triển các ứng dụng robot kết hợp trí tuệ nhân tạo.

Từ các nhu cầu trên, đề tài \textbf{``Nghiên cứu và phát triển robot trợ lý ảo ứng dụng trí tuệ nhân tạo hỗ trợ chăm sóc và giám sát người cao tuổi tại nhà''} xây dựng một nguyên mẫu tích hợp thay vì các chức năng rời rạc. Hệ thống gồm ứng dụng RobotApp trên thiết bị Android gắn với phần thân robot hỗ trợ di chuyển, ứng dụng CaregiverApp cho người chăm sóc, firmware ESP32 điều khiển hai động cơ BLDC ở thân robot và chuỗi mô hình AI được chuyển đổi để suy luận trên thiết bị. Giá trị cốt lõi nằm ở việc kết nối các giải pháp công nghệ thành một hệ thống tương đối hoàn chỉnh hỗ trợ chăm sóc người cao tuổi tại nhà.

\subsection{Mục tiêu đề tài}

\subsubsection{Mục tiêu tổng quát}

Nghiên cứu, thiết kế và xây dựng nguyên mẫu robot trợ lý ảo hỗ trợ chăm sóc người cao tuổi tại nhà. Hệ thống hướng đến khả năng tương tác bằng giọng nói, nhắc nhở theo lịch, giám sát an toàn, phát hiện tình huống nghi té ngã, liên lạc trực tuyến và hỗ trợ người chăm sóc can thiệp từ xa thông qua một nền tảng robot di động sử dụng điện thoại thông minh làm bộ xử lý chính.

\subsubsection{Mục tiêu cụ thể}

\begin{enumerate}
    \item \textbf{Nhận biết và giám sát an toàn:} xây dựng khả năng thu nhận hình ảnh từ môi trường trong nhà, phát hiện người và duy trì đúng đối tượng cần theo dõi khi người hoặc robot di chuyển. Trên cơ sở thông tin tư thế theo thời gian, hệ thống cần nhận biết tình huống nghi té ngã, phân biệt với một số hoạt động thông thường và thực hiện bước xác nhận trước khi phát cảnh báo. Mô hình AI phải được chuẩn bị, chuyển đổi và đánh giá để có thể suy luận trực tiếp trên điện thoại gắn với robot.

    \item \textbf{Tương tác và hỗ trợ chăm sóc:} xây dựng khả năng trò chuyện bằng giọng nói tiếng Việt theo hướng gần gũi, dễ sử dụng đối với người cao tuổi. Hệ thống cần hỗ trợ quản lý và thực hiện lịch nhắc nhở, ghi nhận phản hồi của người dùng, đồng thời kết hợp hội thoại với quy trình xác nhận khi phát hiện dấu hiệu bất thường. Nội dung nhắc nhở và thông tin cá nhân phải được quản lý rõ ràng để người chăm sóc có thể cập nhật khi cần.

    \item \textbf{Liên lạc và điều khiển từ xa:} xây dựng hai ứng dụng dành riêng cho robot và người chăm sóc, cho phép đồng bộ hồ sơ, lịch nhắc nhở, trạng thái giám sát, cảnh báo và lịch sử hoạt động. Khi cần can thiệp, hệ thống phải hỗ trợ cuộc gọi âm thanh, hình ảnh hai chiều và cho phép người chăm sóc điều khiển chuyển động của robot từ xa. Phần thân robot cần tiếp nhận lệnh từ điện thoại và điều khiển hai bánh xe theo yêu cầu di chuyển.

    \item \textbf{An toàn và đánh giá hệ thống:} thiết kế cơ chế dừng robot khi lệnh không hợp lệ, quá thời gian cho phép hoặc đường truyền bị gián đoạn; đồng thời tránh xung đột giữa chế độ tự bám theo và điều khiển thủ công. Dữ liệu cá nhân, lịch sử chăm sóc và thông tin liên lạc cần được giới hạn theo đúng người sử dụng. Cuối cùng, các chức năng chính, tốc độ xử lý, độ ổn định của kết nối, chất lượng phát hiện té ngã và khả năng dừng an toàn phải được kiểm thử trên nguyên mẫu để xác định mức độ đáp ứng và những giới hạn còn tồn tại.
\end{enumerate}

\subsection{Đối tượng và phạm vi nghiên cứu}

\subsubsection{Đối tượng nghiên cứu}

Đối tượng kỹ thuật của đề tài gồm nền tảng robot sử dụng điện thoại thông minh làm bộ xử lý chính; hai ứng dụng Android dành cho robot và người chăm sóc; thị giác máy tính và học sâu trên thiết bị; bài toán theo dõi đối tượng và phân tích chuỗi tư thế; công nghệ chuyển văn bản thành giọng nói (Text-to-Speech -- TTS) và chuyển giọng nói thành văn bản (Speech-to-Text -- STT); mô hình ngôn ngữ lớn; các phương thức đồng bộ dữ liệu và truyền thông thời gian thực; kết nối giữa điện thoại với robot; thiết kế firmware cho hệ thống điều khiển phần thân robot; cùng quy trình chuyển đổi mô hình AI để triển khai trên thiết bị di động.

Đối tượng sử dụng gồm người cao tuổi sinh hoạt tại nhà và người thân/người chăm sóc từ xa. Trong nguyên mẫu, RobotApp và CaregiverApp của một cặp thiết bị đăng nhập cùng tài khoản Firebase; toàn bộ kênh điều khiển, trạng thái AI, signaling và nhật ký thời gian thực được cách ly dưới \path|pair_sessions/{uid}|. Cách tổ chức này tách dữ liệu giữa các tài khoản, nhưng mỗi tài khoản hiện mới đại diện cho một cặp robot--caregiver và chưa mô hình hóa nhiều robot hay nhiều vai trò trong cùng gia đình.

\subsubsection{Phạm vi chức năng}

\begin{itemize}
    \item \textbf{Hỗ trợ tương tác với người cao tuổi:} robot tiếp nhận lời nói, tạo phản hồi và phát lại bằng giọng nói để phục vụ các cuộc trò chuyện hằng ngày. Nội dung tương tác có thể được cá nhân hóa ở mức cơ bản dựa trên hồ sơ và cách xưng hô do người chăm sóc thiết lập.

    \item \textbf{Hỗ trợ nhắc nhở:} người chăm sóc có thể tạo, chỉnh sửa, bật hoặc tắt các lịch nhắc nhở cho những công việc cần thực hiện. Đến thời điểm đã thiết lập, robot phát lời nhắc, tiếp nhận phản hồi của người cao tuổi và lưu lại kết quả để người chăm sóc theo dõi.

    \item \textbf{Giám sát và bám theo người:} robot sử dụng hình ảnh từ camera để phát hiện, lựa chọn và duy trì một người làm đối tượng theo dõi chính. Từ vị trí của người trong khung hình, robot thực hiện các chuyển động tiến, lùi, quay hoặc dừng nhằm giữ khoảng cách quan sát phù hợp trong môi trường trong nhà.

    \item \textbf{Phát hiện và xử lý tình huống nghi té ngã:} hệ thống phân tích sự thay đổi tư thế của người được theo dõi theo thời gian để nhận biết dấu hiệu té ngã. Khi mức nghi ngờ đủ lớn, robot dừng chuyển động và hỏi xác nhận bằng giọng nói; trường hợp người dùng cho biết đang gặp nguy hiểm, trả lời không rõ ràng hoặc không phản hồi sẽ được ghi nhận và thông báo cho người chăm sóc.

    \item \textbf{Liên lạc và can thiệp từ xa:} người chăm sóc có thể nhận thông báo, thực hiện cuộc gọi âm thanh và hình ảnh với người cao tuổi, đồng thời điều khiển chuyển động cơ bản của robot khi cần quan sát hoặc hỗ trợ từ xa. Hệ thống bảo đảm chế độ điều khiển thủ công không hoạt động đồng thời với chế độ tự động bám theo.

    \item \textbf{Quản lý thông tin chăm sóc:} ứng dụng của người chăm sóc hỗ trợ quản lý hồ sơ, lịch nhắc nhở và trạng thái của các chức năng giám sát; đồng thời cung cấp lịch sử thực hiện lời nhắc, cảnh báo té ngã và cuộc gọi để phục vụ việc theo dõi sau sự kiện.
\end{itemize}

\subsubsection{Giới hạn nghiên cứu}

Đề tài được thực hiện ở mức nguyên mẫu kỹ thuật nhằm kiểm chứng sự phối hợp giữa các chức năng hỗ trợ chăm sóc, không hướng đến chẩn đoán bệnh hoặc thay thế quyết định của nhân viên y tế và người chăm sóc. Phạm vi nghiên cứu được giới hạn như sau:

\begin{itemize}
    \item \textbf{Đối tượng và dữ liệu đánh giá:} các tình huống té ngã được sử dụng trong quá trình phát triển và kiểm thử chủ yếu là dữ liệu công khai hoặc kịch bản dựng lại có kiểm soát. Đề tài chưa tiến hành thử nghiệm dài hạn trên người cao tuổi trong điều kiện sinh hoạt thực tế và chưa thực hiện đánh giá lâm sàng.

    \item \textbf{Không gian hoạt động:} robot được định hướng sử dụng trong nhà, trên bề mặt tương đối bằng phẳng và trong điều kiện camera có thể quan sát người dùng. Các môi trường ngoài trời, cầu thang, ánh sáng quá yếu, vật cản phức tạp hoặc rung động mạnh của camera chưa được khảo sát đầy đủ.

    \item \textbf{Khả năng giám sát:} hệ thống ưu tiên theo dõi một người mục tiêu tại một thời điểm và phát hiện tình huống nghi té ngã từ hình ảnh. Kết quả có thể bị ảnh hưởng khi cơ thể bị che khuất, người đi ra ngoài vùng quan sát hoặc xuất hiện nhiều người có đặc điểm gần giống nhau. Hệ thống không đo dấu hiệu sinh tồn, không dự báo nguy cơ té ngã và không bảo đảm phát hiện được mọi dạng sự cố.

    \item \textbf{Khả năng di chuyển:} robot chỉ thực hiện tự bám theo và các lệnh di chuyển cơ bản. Các bài toán lập bản đồ, định vị trong toàn bộ ngôi nhà, lập kế hoạch đường đi và tránh vật cản chủ động chưa thuộc phạm vi nghiên cứu.

    \item \textbf{Kết nối và quy mô sử dụng:} các chức năng đồng bộ, hội thoại trực tuyến, thông báo và gọi từ xa cần có kết nối mạng. Nguyên mẫu được xây dựng cho một robot và một tài khoản chăm sóc; việc quản lý đồng thời nhiều robot, nhiều người cao tuổi hoặc nhiều thành viên chăm sóc chưa được triển khai.

    \item \textbf{Mức độ hoàn thiện sản phẩm:} đề tài chưa bao gồm đánh giá độ bền trong thời gian dài, khả năng chống nước, an toàn điện và cơ khí theo tiêu chuẩn sản phẩm, bảo trì từ xa hoặc chứng nhận dành cho thiết bị y tế. Những nội dung này cần được tiếp tục nghiên cứu trước khi hệ thống có thể được sử dụng như một sản phẩm thương mại.
\end{itemize}

\subsection{Phương pháp nghiên cứu}

Đề tài được thực hiện qua bốn giai đoạn có quan hệ kế tiếp và bổ sung lẫn nhau: khảo cứu tài liệu, phân tích và thiết kế hệ thống, xây dựng nguyên mẫu, sau đó tiến hành thực nghiệm và đánh giá. Kết quả của mỗi giai đoạn được sử dụng để điều chỉnh nội dung ở giai đoạn tiếp theo, nhờ đó hệ thống được hoàn thiện dần theo các mục tiêu đã đặt ra.

\subsubsection{Khảo cứu và phân tích tài liệu}

Phương pháp khảo cứu tài liệu được sử dụng để xác định bối cảnh của bài toán chăm sóc người cao tuổi, các yêu cầu đối với robot trợ lý và những hướng tiếp cận đã được công bố. Nguồn tham khảo gồm báo cáo của các tổ chức có uy tín, bài báo khoa học, tiêu chuẩn kỹ thuật và tài liệu chuyên ngành liên quan đến robot hỗ trợ, thị giác máy tính, theo dõi người, phát hiện té ngã, tương tác bằng giọng nói và hệ thống nhúng. Kết quả khảo cứu là cơ sở để hình thành yêu cầu, lựa chọn hướng giải quyết và xác định phạm vi phù hợp với đồ án.

Đối với bài toán phát hiện té ngã, các nghiên cứu được xem xét theo loại dữ liệu đầu vào và phương pháp nhận biết sự kiện, từ thiết bị đeo, cảm biến môi trường đến phương pháp dựa trên thị giác \reportcite{ref:newaz_fall_review}. Các phương án được phân tích về khả năng ứng dụng trong nhà, mức độ thuận tiện đối với người cao tuổi, yêu cầu phần cứng và khả năng tích hợp với robot di động. Trên cơ sở đó, đề tài tập trung vào hướng phân tích tư thế và chuyển động từ camera gắn trên robot.

\subsubsection{Phân tích yêu cầu và thiết kế hệ thống}

Từ các tình huống chăm sóc tại nhà, nhóm xác định những người tham gia, nhu cầu sử dụng, chức năng cần có, điều kiện vận hành và các trường hợp bất thường có thể xảy ra. Các yêu cầu sau đó được phân nhóm và chuyển thành kiến trúc tổng quan, các khối chức năng cùng những luồng xử lý chính của hệ thống. Quá trình thiết kế chú trọng phân chia rõ trách nhiệm giữa ứng dụng trên robot, ứng dụng của người chăm sóc, khối xử lý AI và phần thân robot; đồng thời xem xét khả năng phối hợp, sử dụng tài nguyên, bảo vệ dữ liệu và dừng an toàn ngay từ đầu.

\subsubsection{Xây dựng và tích hợp nguyên mẫu}

Trên cơ sở thiết kế đã xây dựng, nguyên mẫu được phát triển theo từng thành phần gồm ứng dụng trên robot, ứng dụng của người chăm sóc, khối xử lý AI và chương trình điều khiển phần thân robot. Mỗi thành phần được hoàn thiện và kiểm tra riêng trước khi tích hợp vào hệ thống chung. Quá trình tích hợp được thực hiện lần lượt theo các nhóm chức năng: quản lý hồ sơ và lịch nhắc nhở, hội thoại bằng giọng nói, giám sát và phát hiện té ngã, liên lạc từ xa, điều khiển chuyển động và xử lý các tình huống mất kết nối. Những vấn đề phát sinh trong quá trình chạy thử được ghi nhận để điều chỉnh thiết kế và hoàn thiện nguyên mẫu.

\subsubsection{Thực nghiệm và đánh giá}

Nguyên mẫu được kiểm tra từ từng chức năng riêng lẻ đến các luồng hoạt động hoàn chỉnh trong môi trường trong nhà. Phần xử lý AI được đánh giá về chất lượng nhận biết, khả năng duy trì người mục tiêu và tốc độ xử lý trên điện thoại. Các chức năng tương tác, nhắc nhở, cảnh báo, cuộc gọi và điều khiển từ xa được kiểm tra theo những tình huống sử dụng dự kiến. Phần thân robot được thử nghiệm về hướng chuyển động, khả năng đáp ứng lệnh và trạng thái dừng khi xảy ra bất thường. Ngoài điều kiện vận hành bình thường, các tình huống như mất dấu người, người dùng không phản hồi, kết nối bị gián đoạn hoặc một thành phần ngừng hoạt động cũng được xem xét. Kết quả thực nghiệm được dùng để đối chiếu với mục tiêu đề tài, phân tích hạn chế và đề xuất hướng phát triển tiếp theo.

\subsection{Các công trình và giải pháp liên quan}

Robot hỗ trợ người cao tuổi là một bài toán liên ngành, trong đó khả năng di chuyển, quan sát, tương tác và kết nối với người chăm sóc phải được phối hợp trong cùng một hệ thống. Vì vậy, phần tổng quan tập trung vào bốn nhóm công trình có liên hệ trực tiếp với đề tài: phát hiện té ngã, theo dõi và bám người mục tiêu, hội thoại bằng giọng nói và các nền tảng robot hỗ trợ tại nhà. Mỗi nhóm được xem xét theo ý tưởng chính, cách tổ chức luồng xử lý, ưu điểm, hạn chế và khả năng kế thừa vào nguyên mẫu.

Đối với phát hiện té ngã, các bài tổng quan cung cấp cái nhìn hệ thống về cảm biến và phương pháp đã được sử dụng \reportcite{ref:fall_review}\reportcite{ref:newaz_fall_review}, trong khi những nghiên cứu gần đây cho thấy xu hướng chuyển sang mô hình thị giác nhẹ, khai thác tư thế theo thời gian và xử lý trên thiết bị biên \reportcite{ref:falldetection}\reportcite{ref:hfdmia_pose}\reportcite{ref:yosap_lstm}\reportcite{ref:fallnet}. Việc kết hợp hai nhóm tài liệu này giúp tránh đánh giá một giải pháp chỉ qua độ chính xác được báo cáo, đồng thời làm rõ hơn các yêu cầu về dữ liệu, tốc độ, che khuất và khả năng vận hành trong điều kiện thực tế.

\subsubsection{Các hướng tiếp cận cho bài toán phát hiện té ngã}

\paragraph{Bức tranh chung về các phương pháp phát hiện té ngã.}
Các hệ thống phát hiện té ngã được phát triển từ nhiều nguồn dữ liệu khác nhau, bao gồm gia tốc và vận tốc góc, tín hiệu sinh học như điện tâm đồ (ECG) thu từ thiết bị đeo, âm thanh, ảnh chiều sâu, radar và video RGB. Hai tổng quan gần đây có cách phân nhóm khác nhau nhưng cùng cho thấy hiệu quả của hệ thống phụ thuộc đồng thời vào cảm biến, thuật toán, bộ dữ liệu và phương pháp đánh giá \reportcite{ref:fall_review}\reportcite{ref:newaz_fall_review}. Vì vậy, tỷ lệ xuất hiện của một hướng nghiên cứu không đồng nghĩa với ưu thế tuyệt đối về độ chính xác hoặc khả năng triển khai.

Xét theo cách thu nhận bằng chứng, các phương pháp có thể quy về ba nhóm chính: cảm biến gắn trên cơ thể, cảm biến bố trí trong môi trường và camera quan sát. Điện thoại thông minh, kết nối IoT hoặc điện toán đám mây chủ yếu giữ vai trò xử lý, lưu trữ và chuyển tiếp cảnh báo nên có thể xuất hiện đồng thời trong cả ba nhóm. Cách phân chia theo nguồn dữ liệu phù hợp hơn khi đánh giá một robot hoạt động trong nhà, bởi nó làm rõ người dùng có phải mang thiết bị hay không, phạm vi quan sát, mức độ riêng tư và phần cứng cần bổ sung.

\paragraph{Phương pháp sử dụng thiết bị đeo và tín hiệu cơ thể.}
Các hệ thống thuộc nhóm này thường sử dụng gia tốc kế, con quay hồi chuyển hoặc IMU gắn tại thắt lưng, ngực, cổ tay hay tích hợp sẵn trong điện thoại. Tín hiệu được lấy mẫu, lọc nhiễu và chia thành các cửa sổ thời gian trước khi đưa qua luật ngưỡng hoặc mô hình phân loại để xác nhận sự kiện. Một đại lượng thường được sử dụng là độ lớn vector gia tốc
\begin{equation}
a=\sqrt{a_x^2+a_y^2+a_z^2},
\end{equation}
sau đó đối chiếu diễn biến giảm tải, va chạm, thay đổi góc thân và trạng thái bất động. Bourke, O'Brien và Lyons đã đánh giá một thuật toán ngưỡng sử dụng gia tốc kế ba trục gắn trên cơ thể, qua đó cho thấy khả năng hiện thực bộ phát hiện có chi phí tính toán thấp \reportcite{ref:bourke_threshold}. Trên nền tảng smartphone, Abbate và cộng sự bổ sung bước nhận biết những hoạt động thường ngày dễ bị nhầm với té ngã trước khi tự động gửi yêu cầu trợ giúp \reportcite{ref:abbate_smartphone}. Việc xem xét cả trạng thái trước và sau va chạm giúp giảm sự phụ thuộc vào một đỉnh gia tốc đơn lẻ.

Các phương pháp dựa trên chuyển động có chi phí tính toán thấp, không ghi lại hình ảnh trong không gian sống, ít chịu ảnh hưởng của điều kiện chiếu sáng và vẫn hoạt động khi người dùng ở ngoài tầm quan sát của camera. Khi luật ngưỡng được thay bằng SVM, random forest, CNN một chiều hoặc LSTM/GRU, mô hình có khả năng học quan hệ giữa nhiều pha chuyển động. Bộ dữ liệu SisFall, được thu từ hai gia tốc kế và một con quay đặt tại thắt lưng, bao gồm nhiều hoạt động thường ngày và kiểu ngã khác nhau \reportcite{ref:sisfall}. Kết quả đánh giá thay đổi đáng kể giữa các nhóm tuổi cho thấy mô hình huấn luyện chủ yếu bằng những cú ngã dàn dựng của người trẻ có thể suy giảm khả năng khái quát khi áp dụng cho người cao tuổi.

Bên cạnh tín hiệu chuyển động, một số nghiên cứu khai thác ECG để ghi nhận hoạt động điện của tim và những biến đổi sinh lý có thể xuất hiện trước, trong hoặc sau sự cố. Butt và cộng sự sử dụng dữ liệu ECG của các trường hợp ngã và không ngã, sau đó áp dụng học chuyển giao với AlexNet và GoogLeNet để phân loại; mô hình tốt nhất đạt độ chính xác 98,44\% trong điều kiện thực nghiệm của nghiên cứu \reportcite{ref:ecg_fall_transfer}. Kết quả này cho thấy tín hiệu sinh học có thể cung cấp thêm bằng chứng mà gia tốc hoặc hình ảnh không phản ánh trực tiếp, đồng thời mở ra khả năng kết hợp phát hiện sự cố với theo dõi trạng thái tim mạch.

Tuy nhiên, biến đổi ECG không phải lúc nào cũng bắt nguồn từ té ngã và bản thân tín hiệu này không cho biết vị trí hay tư thế của người trong không gian. Chất lượng đo còn phụ thuộc vị trí điện cực, mức tiếp xúc với da và nhiễu do chuyển động. Các phương pháp dùng thiết bị đeo nói chung cũng phụ thuộc vào việc người dùng đeo đúng vị trí, duy trì đủ pin và không bỏ quên thiết bị. Với cảm biến quán tính, những hoạt động như chạy, nhảy hoặc ngồi mạnh có thể tạo xung gia tốc gần giống va chạm, trong khi một cú trượt chậm lại có khả năng không vượt ngưỡng phát hiện. Vì vậy, ECG và IMU phù hợp hơn khi được sử dụng như nguồn thông tin bổ trợ. Trong phạm vi đề tài, camera vẫn là lựa chọn chính vì vừa cung cấp tư thế, vừa xác định vị trí người để phục vụ theo dõi và điều khiển robot; tín hiệu sinh học là một hướng mở rộng cho chức năng giám sát sức khỏe sau này.

\paragraph{Phương pháp dựa trên cảm biến môi trường.}
Thay vì gắn cảm biến lên cơ thể, nhóm phương pháp này bố trí thiết bị thu nhận trong không gian cần giám sát. Acoustic-FADE sử dụng mảng microphone và SRP-PHAT để định vị nguồn âm, beamforming để tăng cường tín hiệu, MFCC để biểu diễn đặc trưng phổ, đồng thời khai thác độ cao nguồn âm nhằm hỗ trợ phân biệt tiếng ngã với các âm thanh khác \reportcite{ref:acoustic_fade}. Trên tập thử nghiệm gồm 120 âm thanh ngã và 120 âm thanh không ngã được mô phỏng, hệ thống đạt độ nhạy 100\% tại độ đặc hiệu 97\%. Mặc dù cho kết quả cao trong điều kiện thí nghiệm, tín hiệu âm thanh dễ thay đổi theo tiếng nền, độ vang, vật liệu sàn và âm thanh do đồ vật rơi; dữ liệu dàn dựng cũng chưa phản ánh đầy đủ sự đa dạng của tai nạn tự nhiên.

Camera chiều sâu và hệ stereo bổ sung quan hệ hình học ba chiều giữa cơ thể với mặt sàn. Stone và Skubic xây dựng một quy trình hai giai đoạn trên Kinect: trước hết xác định các đoạn chuỗi trong đó người ở gần sàn, sau đó dùng tập hợp cây quyết định để ước lượng khả năng đã xảy ra cú ngã \reportcite{ref:kinect_fall}. Công trình được đánh giá trong thời gian dài tại 13 căn hộ, ghi nhận 454 lần ngã nhưng chỉ có chín trường hợp diễn ra tự nhiên. Sự chênh lệch này phản ánh khó khăn trong việc thu thập dữ liệu té ngã thực tế có nhãn đáng tin cậy. Solbach và Tsotsos kết hợp CNN ước lượng tư thế với camera stereo để khôi phục pose 3D và mặt phẳng sàn \reportcite{ref:stereo_fall}. Thông tin hình học giúp giảm ảnh hưởng của màu sắc và nền cảnh, nhưng yêu cầu phần cứng, phạm vi đo và quy trình hiệu chuẩn riêng.

Radar mmWave khai thác khoảng cách, vận tốc Doppler hoặc point cloud mà không tạo ra hình ảnh có thể nhận dạng khuôn mặt, đồng thời vẫn hoạt động trong điều kiện thiếu sáng. Hệ thống mmFall sử dụng point cloud 4D và hybrid variational RNN autoencoder để học phân bố của các hoạt động bình thường, sau đó kết hợp điểm bất thường với sự hạ thấp của tâm cơ thể \reportcite{ref:mmfall}. Trong 50 cú ngã thử nghiệm, phương pháp phát hiện được 98\% và tạo hai báo động giả. Tuy nhiên, độ thưa của point cloud, phản xạ đa đường và cách bố trí phòng có thể làm thay đổi đáng kể tín hiệu. Để khắc phục giới hạn của từng cảm biến riêng lẻ, Kwolek và Kepski sử dụng IMU để kích hoạt một sự kiện tiềm năng, rồi xác minh bằng ảnh chiều sâu, đặc trưng hình dạng, khoảng cách tới sàn và SVM \reportcite{ref:kwolek_fusion}. Hợp nhất cảm biến giúp các nguồn dữ liệu bù trừ điểm mù cho nhau, nhưng kéo theo chi phí phần cứng, đồng bộ thời gian, hiệu chuẩn và bảo trì cao hơn.

\begin{table}[H]
\centering
\caption{So sánh bằng chứng và giới hạn của các hướng phát hiện té ngã tiêu biểu}
\label{tab:fall_sensing_comparison}
\small
\begin{tabularx}{\textwidth}{|p{2.6cm}|X|X|X|}
\hline
\textbf{Hướng} & \textbf{Bằng chứng chính} & \textbf{Ưu điểm} & \textbf{Giới hạn khi triển khai} \\
\hline
IMU/thiết bị đeo & Gia tốc, vận tốc góc, góc thân, bất động & Nhẹ, riêng tư, không phụ thuộc camera & Phải đeo/sạc đúng cách; nhạy vị trí và người dùng \\
\hline
Tín hiệu sinh học & ECG, nhịp tim và biến đổi sinh lý & Bổ sung thông tin về trạng thái cơ thể và nguy cơ liên quan & Cần tiếp xúc cơ thể; nhạy nhiễu chuyển động và không cung cấp vị trí người \\
\hline
Âm thanh & Va chạm, MFCC, vị trí và độ cao nguồn âm & Không đeo, không lưu ảnh & Nhạy tiếng nền, độ vang, đồ vật rơi và vật liệu sàn \\
\hline
Depth/stereo & Silhouette/pose 3D, tâm người, khoảng cách tới sàn & Hình học ba chiều rõ, giảm phụ thuộc màu sắc & Cần phần cứng và hiệu chuẩn; vùng quan sát cố định \\
\hline
Radar mmWave & Doppler, point cloud, quỹ đạo tâm cơ thể & Hoạt động thiếu sáng, bảo vệ riêng tư tốt hơn RGB & Dữ liệu thưa, đa đường, chi phí cảm biến \\
\hline
RGB một khung hình & Hộp người/ngã và đặc trưng hình ảnh & Nhanh, định vị trực tiếp, tận dụng camera sẵn có & Dễ nhầm nằm/cúi; nhạy ánh sáng, góc nhìn và che khuất \\
\hline
RGB pose--thời gian & Keypoint, quỹ đạo khớp, thay đổi tư thế theo chuỗi & Mô tả được diễn biến cú ngã, phù hợp xử lý biên & Lỗi pose/tracking lan sang LSTM; cần dữ liệu chuỗi đa dạng \\
\hline
Đa cảm biến & Quyết định từ hai hoặc nhiều nguồn & Bù điểm mù và tăng khả năng xác minh & Tăng chi phí, đồng bộ, hiệu chuẩn và bảo trì \\
\hline
\end{tabularx}
\end{table}

\paragraph{Phương pháp sử dụng camera và thị giác máy tính.}
Camera cung cấp nhiều thông tin hơn một tín hiệu chuyển động đơn lẻ: từ video, hệ thống có thể xác định vị trí người, hình dạng hộp bao, các điểm khớp, tốc độ hạ thấp cơ thể và trạng thái sau va chạm. Khung hình gắn với sự kiện còn giúp người chăm sóc quan sát lại tình huống khi cần thiết. Sự phổ biến của camera tích hợp và bộ xử lý di động làm cho hướng tiếp cận này ngày càng phù hợp với xử lý tại biên. Tuy nhiên, chất lượng nhận biết vẫn chịu ảnh hưởng của ánh sáng, góc nhìn, che khuất và chuyển động nền; dữ liệu hình ảnh cũng đòi hỏi nguyên tắc bảo vệ quyền riêng tư rõ ràng.

Các phương pháp thị giác ban đầu thường dùng silhouette, tỷ lệ chiều cao--chiều rộng, hướng trục chính hoặc lịch sử chuyển động. Hướng optical flow biểu diễn chuyển động giữa các khung hình; Núñez-Marcos, Azkune và Arganda-Carreras xếp chồng optical flow rồi dùng CNN dựa trên VGG-16 để phân lớp đoạn video \reportcite{ref:opticalflow_fall}. Cách này nắm được động học tốt hơn ảnh tĩnh, nhưng bước tính dòng quang học tốn tài nguyên và đặc biệt nhạy với chuyển động camera của robot.

Với detector một giai đoạn, ảnh RGB đi qua mạng YOLO, ngưỡng độ tin cậy và NMS để trả về hộp người ngã. Công trình ``Visionary Vigilance'' xây dựng DiverseFALL10500 với khoảng 10.500 ảnh đa dạng và tối ưu YOLOv8s bằng Focus cùng CBAM \reportcite{ref:falldetection}. Detector ảnh có tốc độ tốt và định vị trực tiếp sự cố, nhưng một tư thế nằm trong ảnh đơn không chứng minh đã có chuyển tiếp té ngã: người dùng có thể nằm nghỉ, cúi nhặt đồ, tập thể dục hoặc bị quan sát ở góc bất lợi. Vì vậy, accuracy theo khung hình cao vẫn có thể dẫn đến nhiều cảnh báo giả khi camera hoạt động hàng giờ.

Ước lượng pose chuyển người thành tập keypoint $(x,y,c)$, trong đó $c$ là độ tin cậy. Sau khi chuẩn hóa theo hộp người hoặc kích thước cơ thể, hệ thống có thể tạo các đặc trưng như góc thân, tỷ lệ chiều cao--chiều rộng, độ cao tương đối và vận tốc của hông, vai, gối. So với pixel RGB, skeleton giảm phụ thuộc vào màu quần áo và nền, đồng thời làm rõ diễn biến động học. Tuy nhiên, che khuất và góc nhìn xấu có thể làm mất keypoint; bởi vậy pose phải được gắn với đúng track trước khi ghép thành chuỗi. Juraev và cộng sự cũng chỉ ra các khó khăn về thiếu dữ liệu ngã thực, thay đổi góc camera/ánh sáng và che khuất; dữ liệu tổng hợp có thể hỗ trợ, nhưng lựa chọn mô hình thời gian vẫn phải được kiểm chứng trên miền triển khai \reportcite{ref:pose_synthetic}.

\paragraph{Kết hợp tư thế với thông tin thời gian.}
Hạn chế của quyết định dựa trên từng khung hình tạo động lực kết hợp bộ ước lượng tư thế với mô hình chuỗi. Trong kiến trúc này, YOLO-pose xác định vị trí người và các khớp ở mỗi thời điểm, còn LSTM khai thác sự phụ thuộc giữa các khung liên tiếp để phân biệt một quá trình té ngã với tư thế nằm hoặc cúi người thông thường. Luồng xử lý tổng quát gồm
\begin{center}
\small
\textit{Video RGB $\rightarrow$ ước lượng tư thế và gán ID $\rightarrow$ chuỗi keypoint đã chuẩn hóa}\\
\textit{$\rightarrow$ mô hình thời gian $\rightarrow$ xác nhận sự kiện $\rightarrow$ phát cảnh báo.}
\end{center}
Nhờ duy trì trạng thái theo thời gian, LSTM có thể học sự nối tiếp giữa mất thăng bằng, hạ thấp nhanh, va chạm và trạng thái nằm bất thường thay vì chỉ nhận dạng một dáng người đơn lẻ \reportcite{ref:lstm}. YOLO-pose và LSTM vì thế đảm nhiệm hai vai trò bổ sung: mô hình đầu trích xuất biểu diễn không gian, còn mô hình sau diễn giải sự biến đổi của biểu diễn đó theo thời gian.

FallNet là một hiện thực tiêu biểu của kiến trúc này. Mô hình được phát triển từ YOLOv8n-pose, trong đó khối C2F được thay bằng DyC2F nhằm cải thiện khả năng trích xuất đặc trưng động và giảm chi phí tính toán; quỹ đạo keypoint sau đó được đưa vào LSTM để nhận biết sớm xu hướng ngã \reportcite{ref:fallnet}. Nghiên cứu báo cáo độ chính xác 97\%, precision 92\%, recall 97\% và F1-score 94\% với 3,4 triệu tham số cùng 6,4 GFLOPs. Khi triển khai trên Raspberry Pi 3B+, nguyên mẫu đạt 13 FPS với độ trễ xử lý tại thiết bị khoảng 100 ms. Việc báo cáo đồng thời chất lượng phân loại, quy mô mô hình và tốc độ trên phần cứng thực làm cho công trình có giá trị tham khảo trực tiếp đối với các hệ thống AI biên.

Vai trò của bước theo dõi được thể hiện rõ hơn trong YOSAP-LSTM. Hệ thống dùng YOLOv8 để phát hiện người, SORT để duy trì định danh, AlphaPose để trích xuất skeleton và mạng 1D CNN--LSTM để phân loại chuỗi \reportcite{ref:yosap_lstm}. Trên bộ dữ liệu công trường của nghiên cứu, mô hình đạt độ chính xác 98,66\%, độ nhạy 97,32\%, độ đặc hiệu 99,10\% và tốc độ 6,44 FPS trên NVIDIA Jetson Xavier NX. Khả năng xử lý nhiều người và các mức che khuất khác nhau cho thấy tracking là mắt xích cần thiết trước khi hình thành chuỗi pose. Tuy nhiên, miền công trường, cấu hình Jetson và dữ liệu chuyên biệt khác đáng kể so với môi trường nhà ở trên smartphone; do đó kết quả 98,66\% không thể được xem là hiệu năng kỳ vọng của đồ án.

Ở một biến thể khác, Liu và cộng sự đề xuất HFDMIA-Pose, cải tiến YOLOv8s bằng SPD-Conv, bổ sung đầu phát hiện vật thể nhỏ và sử dụng BCIoU; kết quả của detector được chuyển sang AlphaPose và kết hợp với các đặc trưng tức thời cũng như trạng thái để xác định cú ngã \reportcite{ref:hfdmia_pose}. Trên Le2i và bộ MPFDD nhiều người, phương pháp cải thiện tương đối so với AlphaPose 4,30\% về độ chính xác, 4,57\% về F1 và 37,50\% về FPS. Mặc dù không sử dụng LSTM, nghiên cứu cho thấy chất lượng phát hiện người và keypoint quyết định trực tiếp chất lượng của tầng nhận dạng hành động phía sau; đồng thời, bằng chứng tức thời chỉ trở nên đáng tin cậy khi được đặt trong một trạng thái kéo dài theo thời gian.

Bảng~\ref{tab:recent_vision_fall} tóm tắt một số công trình thị giác tiêu biểu được công bố gần đây. Các nghiên cứu này không hoàn toàn cùng bài toán, nhưng thể hiện rõ ba hướng cải tiến đang được quan tâm: tăng tính đa dạng của dữ liệu, khai thác tư thế theo thời gian và giảm chi phí tính toán để đưa mô hình tới thiết bị biên.

\begin{table}[H]
\centering
\caption{Một số nghiên cứu thị giác gần đây về phát hiện té ngã}
\label{tab:recent_vision_fall}
\small
\begin{tabularx}{\textwidth}{|p{2.8cm}|p{4.4cm}|X|}
\hline
\textbf{Công trình} & \textbf{Hướng tiếp cận} & \textbf{Kết quả và điểm cần lưu ý} \\
\hline
Visionary Vigilance (2024) \reportcite{ref:falldetection} & Cải tiến YOLOv8s và xây dựng DiverseFALL10500 & Bộ dữ liệu có quy mô và bối cảnh đa dạng; phát hiện trực tiếp trên ảnh nhanh nhưng chưa mô tả đầy đủ diễn biến của sự kiện. \\
\hline
HFDMIA-Pose (2025) \reportcite{ref:hfdmia_pose} & YOLOv8s cải tiến, AlphaPose và đặc trưng tức thời--trạng thái & Hướng đến nhiều người và thiết bị biên; hiệu quả vẫn phụ thuộc chất lượng pose và mức che khuất. \\
\hline
YOSAP-LSTM (2025) \reportcite{ref:yosap_lstm} & YOLOv8, SORT, AlphaPose và 1D CNN--LSTM & Đạt 98,66\% accuracy và 6,44 FPS trên Jetson Xavier NX; miền dữ liệu công trường khác với môi trường nhà ở. \\
\hline
FallNet (2026) \reportcite{ref:fallnet} & YOLOv8n-pose cải tiến kết hợp LSTM & Đạt 97\% accuracy với 3,4 triệu tham số; cho thấy tiềm năng của pose--LSTM trên thiết bị có tài nguyên hạn chế. \\
\hline
\end{tabularx}
\end{table}

Các con số trên không nên đặt cạnh nhau như một bảng xếp hạng. FallNet, YOSAP-LSTM và HFDMIA-Pose dùng miền dữ liệu, định nghĩa nhãn, cách chia tập, phần cứng và đơn vị đo khác nhau. Accuracy còn có thể che giấu mất cân bằng lớp; trong chăm sóc dài hạn, recall của cú ngã, precision, F1, số báo giả mỗi giờ, độ trễ phát hiện và khả năng tổng quát sang người/căn phòng chưa thấy đều quan trọng. Đánh giá đáng tin cậy cần tách train/test theo người hoặc video, không chia ngẫu nhiên các khung gần nhau của cùng một cảnh, và phải thử bổ sung trên camera thực của robot. CAUCAFall là một nguồn dữ liệu hữu ích do có thay đổi khoảng cách, hướng ngã và ánh sáng trong môi trường ít kiểm soát, nhưng các cú ngã vẫn chủ yếu phải dàn dựng vì lý do đạo đức và an toàn \reportcite{ref:caucafall}.

\paragraph{Hướng tiếp cận được lựa chọn cho đề tài.}
Từ các kết quả đã phân tích, hướng kết hợp pose và mô hình thời gian phù hợp với nguyên mẫu robot chăm sóc trên cả phương diện chức năng lẫn khả năng triển khai. Camera trên điện thoại cho phép quan sát không tiếp xúc, không yêu cầu người cao tuổi luôn mang thêm cảm biến. Hộp người và keypoint còn có thể được sử dụng chung cho chức năng theo dõi mục tiêu, trong khi chuỗi tư thế cung cấp thêm bằng chứng để phân biệt cú ngã với hành động nằm hoặc cúi người có chủ đích.

Nhược điểm của kiến trúc này là sai số có thể truyền qua nhiều giai đoạn: bỏ sót người làm gián đoạn track, track không ổn định làm trộn chuỗi tư thế và keypoint sai tiếp tục ảnh hưởng đến kết quả phân loại. Do đó, dự đoán của mô hình chỉ nên được xem là một nguồn bằng chứng và cần được kiểm tra theo nhiều khung hình trước khi tạo cảnh báo. Đề tài lựa chọn hướng YOLO-pose kết hợp LSTM vì đạt được sự cân bằng giữa thông tin tư thế, chi phí tính toán và khả năng tích hợp với robot di động; các tham số và cơ chế hậu kiểm cụ thể được trình bày tại phần thiết kế và hiện thực hệ thống.

\iffalse

\paragraph{Phát hiện bằng thiết bị đeo và cảm biến môi trường.}
Thiết bị đeo thường khai thác gia tốc kế, con quay hồi chuyển hoặc tổ hợp cảm biến quán tính (IMU) để tìm biến đổi đột ngột của gia tốc, góc cơ thể và trạng thái bất động sau va chạm. Luồng điển hình là \textit{lấy mẫu cảm biến $\rightarrow$ lọc nhiễu $\rightarrow$ chia cửa sổ thời gian $\rightarrow$ trích đặc trưng $\rightarrow$ ngưỡng hoặc bộ phân loại $\rightarrow$ xác nhận và gửi cảnh báo}. Ưu điểm là không ghi hình không gian sống, tín hiệu gắn trực tiếp với người đeo và có thể hoạt động khi người dùng rời khỏi vùng quan sát của camera. Đổi lại, kết quả phụ thuộc vị trí gắn, hướng cảm biến, tần số lấy mẫu và việc người cao tuổi có duy trì đeo/sạc thiết bị hay không.

\paragraph{Luật ngưỡng trên tín hiệu quán tính.}
Nhóm phương pháp đơn giản nhất tính độ lớn vector gia tốc $a=\sqrt{a_x^2+a_y^2+a_z^2}$, tìm đỉnh va chạm và kiểm tra tư thế hoặc mức vận động sau đỉnh. Bourke, O'Brien và Lyons đã đánh giá thuật toán ngưỡng trên gia tốc kế ba trục gắn cơ thể, qua đó hình thành một mốc tham chiếu sớm cho thiết kế fall detector ít tính toán \reportcite{ref:bourke_threshold}. Cách tiếp cận này phù hợp vi điều khiển hoặc điện thoại vì chỉ cần bộ lọc, cửa sổ trượt và các phép so sánh. Tuy nhiên, ngưỡng đặt cho một nhóm người hoặc vị trí đeo không bảo đảm chuyển sang nhóm khác; chạy, nhảy, ngồi mạnh và đặt thiết bị xuống bàn có thể tạo đỉnh gia tốc tương tự.

Một cải tiến quan trọng là không quyết định chỉ từ đỉnh va chạm mà kiểm tra chuỗi pha: rơi tự do hoặc giảm tải, va chạm, thay đổi góc thân và bất động sau ngã. Hệ thống smartphone của Abbate và cộng sự theo dõi chuyển động, nhận biết ngã rồi tự động yêu cầu trợ giúp; nhóm tác giả bổ sung nhận dạng các hoạt động thường ngày dễ bị nhầm như ngồi xuống sofa hoặc nằm lên giường để giảm báo động giả \reportcite{ref:abbate_smartphone}. Ý tưởng này tương đồng với hậu kiểm của đồ án: một bằng chứng mạnh tức thời cần được xác nhận bằng trạng thái trước và sau sự kiện.

\paragraph{Học máy và học sâu trên dữ liệu IMU.}
Thay vì chọn ngưỡng thủ công, tín hiệu gia tốc và vận tốc góc có thể được chia thành các cửa sổ rồi đưa vào SVM, random forest, CNN, LSTM hoặc mô hình kết hợp. Đặc trưng đầu vào có thể là giá trị thống kê theo trục, năng lượng, độ nghiêng, miền tần số hoặc trực tiếp chuỗi thô. SisFall là bộ dữ liệu có ảnh hưởng trong hướng này: thiết bị đeo ở thắt lưng gồm hai gia tốc kế và một con quay; 23 người trẻ thực hiện 19 hoạt động thường ngày và 15 kiểu ngã, 14 người trên 62 tuổi thực hiện 15 hoạt động thường ngày, còn một người 60 tuổi thực hiện cả hai nhóm \reportcite{ref:sisfall}. Đáng chú ý, nghiên cứu cho thấy hiệu năng giảm khi kiểm tra với nhóm cao tuổi, cảnh báo nguy cơ mô hình học đặc điểm của người trẻ thực hiện cú ngã dàn dựng thay vì đặc điểm có khả năng khái quát.

Với học sâu, CNN một chiều có thể học mẫu cục bộ của tín hiệu, còn LSTM/GRU mô hình hóa thứ tự các pha của cú ngã. Các mô hình này giảm phụ thuộc vào đặc trưng thiết kế tay nhưng cần nhiều dữ liệu, dễ học lệch từ vị trí cảm biến và khó giải thích hơn luật ngưỡng. Một hướng khác là \emph{pre-impact detection}: phát hiện giai đoạn mất thăng bằng trước khi cơ thể chạm sàn để kích hoạt thiết bị bảo vệ. Bài toán này đòi hỏi độ trễ rất thấp và đánh đổi khắt khe giữa phát hiện sớm với báo giả; nó khác với mục tiêu của đồ án là xác nhận sự cố và liên lạc sau ngã.

\paragraph{Cảm biến âm thanh và rung động.}
Cảm biến môi trường loại bỏ yêu cầu người dùng mang thiết bị. Hệ acoustic-FADE của Li, Ho và Popescu dùng mảng microphone để định vị nguồn âm bằng SRP-PHAT, tăng cường tín hiệu bằng beamforming, trích MFCC và phân lớp fall/non-fall; độ cao nguồn âm được dùng như bằng chứng loại bỏ âm thanh phát ra phía trên sàn \reportcite{ref:acoustic_fade}. Trên bộ dữ liệu 120 âm thanh ngã và 120 không ngã được mô phỏng, công trình báo cáo sensitivity 100\% tại specificity 97\%. Kết quả cho thấy âm thanh có thể hỗ trợ tốt mà không tạo ảnh người dùng, nhưng tiếng đồ vật rơi, tiếng nền, độ vang phòng và sự khác nhau của vật liệu sàn làm thay đổi mạnh miền tín hiệu. Ngoài ra, dữ liệu mô phỏng chưa phản ánh đầy đủ âm thanh của một tai nạn tự nhiên.

\paragraph{Camera chiều sâu và cảm biến ba chiều.}
Camera depth biểu diễn khoảng cách thay cho màu sắc chi tiết, nhờ đó có thể theo dõi chiều cao cơ thể và giảm một phần lo ngại riêng tư. Stone và Skubic xây dựng hệ hai giai đoạn trên Microsoft Kinect: giai đoạn đầu suy ra trạng thái thẳng đứng theo từng ảnh depth và tách các đoạn ``ở trên sàn'' từ chuỗi theo dõi; giai đoạn sau dùng tập hợp cây quyết định để ước lượng khả năng một cú ngã xảy ra trước trạng thái đó \reportcite{ref:kinect_fall}. Điểm mạnh của nghiên cứu là đánh giá trên tổng cộng chín năm dữ liệu liên tục tại 13 căn hộ với 454 lần ngã, nhưng chỉ chín trường hợp là ngã tự nhiên; con số này đồng thời cho thấy mức độ khó của việc thu thập ground truth thực địa.

Solbach và Tsotsos đi theo hướng stereo: CNN ước lượng pose, cặp camera khôi phục tư thế 3D và mặt phẳng sàn, sau đó mô-đun suy luận kiểm tra quan hệ cơ thể--mặt sàn để nhận biết người đã ngã \reportcite{ref:stereo_fall}. Hình học 3D giúp giảm sự phụ thuộc vào góc nhìn và tỷ lệ hộp 2D, nhưng camera stereo/depth có chi phí, giới hạn khoảng đo và yêu cầu hiệu chuẩn; đây là lý do đồ án ưu tiên camera RGB sẵn có trên smartphone.

Kwolek và Kepski kết hợp hai loại bằng chứng theo cơ chế kích hoạt--xác minh: gia tốc kế ba trục phát hiện cú ngã tiềm năng; ảnh depth sau đó được xử lý để tách người, tính đặc trưng hình dạng và khoảng cách tới sàn, rồi SVM xác thực cảnh báo \reportcite{ref:kwolek_fusion}. Nếu mục tiêu nằm ngoài trường nhìn, hệ vẫn có thể dùng IMU; nếu người dùng không mang cảm biến, nhánh depth vẫn hoạt động. Đây là ví dụ rõ về sensor fusion nhằm bù điểm mù lẫn nhau. Với đồ án, nguyên tắc có thể được vận dụng ở mức phần mềm bằng cách kết hợp xác suất LSTM, diễn biến pose, độ tin cậy tracking và phản hồi giọng nói, dù phần cứng hiện tại chỉ có camera RGB.

\paragraph{Radar mmWave và phát hiện bất thường.}
Radar đo khoảng cách, vận tốc Doppler hoặc point cloud mà không ghi lại hình ảnh nhận dạng khuôn mặt, hoạt động trong bóng tối và phù hợp khu vực nhạy cảm. mmFall của Jin, Sengupta và Cao thu point cloud 4D cùng tâm cơ thể, sau đó dùng hybrid variational RNN autoencoder học phân bố hoạt động thường ngày \reportcite{ref:mmfall}. Khi suy luận, hệ thống kết hợp đỉnh điểm bất thường với sự giảm độ cao tâm cơ thể để phát hiện ngã. Mô hình chỉ học dữ liệu bình thường giúp giảm nhu cầu thu thập cú ngã thật; thử nghiệm của công trình phát hiện 98\% trong 50 cú ngã với hai báo động giả. Dù vậy, point cloud radar thưa và ngẫu nhiên, chi phí cảm biến cao hơn camera điện thoại, đồng thời phản xạ đa đường phụ thuộc bố trí phòng. Radar vì thế thích hợp như kênh bổ sung hoặc cho khu vực cần riêng tư hơn là thay thế trực tiếp giải pháp smartphone của đồ án.

Các nghiên cứu trên cho thấy ``không đeo'' không đồng nghĩa với ``không có ràng buộc'': microphone phụ thuộc âm học phòng; depth/stereo phụ thuộc tầm nhìn và hiệu chuẩn; radar phụ thuộc phản xạ; cảm biến sàn chỉ bao phủ vùng lắp đặt. Tổng quan của Alam và cộng sự cho thấy không có một loại cảm biến duy nhất giải quyết đồng thời độ bao phủ, riêng tư, chi phí và độ chính xác trong mọi bối cảnh \reportcite{ref:fall_review}; Newaz và Hanada bổ sung góc nhìn hệ thống về kết nối IoT/cloud, bảo vệ dữ liệu và lựa chọn phương pháp kiểm định \reportcite{ref:newaz_fall_review}. Bảng \ref{tab:fall_sensing_comparison} hệ thống hóa các lựa chọn.

\begin{table}[htbp]
\centering
\caption{So sánh các hướng cảm biến cho phát hiện té ngã}
\label{tab:fall_sensing_comparison}
\small
\begin{tabularx}{\textwidth}{|p{2.3cm}|X|X|X|}
\hline
\textbf{Hướng} & \textbf{Flow/đặc trưng chính} & \textbf{Ưu điểm} & \textbf{Giới hạn và liên hệ đề tài} \\
\hline
IMU ngưỡng & Độ lớn gia tốc, góc thân, bất động sau va chạm & Nhẹ, rẻ, chạy cục bộ & Phải đeo và hiệu chỉnh; là tham chiếu cho hậu kiểm nhiều pha \\
\hline
IMU học máy & Cửa sổ gia tốc/con quay đưa vào SVM, CNN hoặc LSTM & Học được mẫu phức tạp theo thời gian & Lệch miền theo người/vị trí đeo; đồ án dùng LSTM nhưng trên pose \\
\hline
Âm thanh/rung & Định vị nguồn, beamforming, MFCC hoặc đặc trưng rung & Không đeo, không ghi hình & Nhạy tiếng nền và vật liệu phòng; có thể là cảm biến bổ sung \\
\hline
Depth/stereo & Silhouette/pose 3D, chiều cao tâm người, khoảng cách tới sàn & Có thông tin hình học ba chiều & Cần phần cứng và hiệu chuẩn riêng; đồ án tận dụng RGB smartphone \\
\hline
Radar mmWave & Doppler/point cloud, RNN hoặc phát hiện bất thường & Hoạt động thiếu sáng, bảo vệ riêng tư tốt hơn RGB & Dữ liệu thưa, phản xạ đa đường, tăng chi phí \\
\hline
Camera RGB & Detector hoặc pose kết hợp mô hình thời gian & Có sẵn, giàu thông tin, phù hợp AI biên & Nhạy ánh sáng/che khuất và cần quản trị dữ liệu hình ảnh \\
\hline
\end{tabularx}
\end{table}

\paragraph{Phát hiện trực tiếp trên ảnh bằng mô hình detector.}
Một hướng phổ biến là gán nhãn người ở trạng thái bình thường và người đã ngã, sau đó huấn luyện bộ phát hiện đối tượng một giai đoạn như YOLO. Pipeline cơ bản gồm \textit{khung hình RGB $\rightarrow$ đổi kích thước/chuẩn hóa $\rightarrow$ mạng detector $\rightarrow$ lọc theo độ tin cậy và NMS $\rightarrow$ hộp ``fall'' $\rightarrow$ cảnh báo}. Cách này có ưu thế về tốc độ và cho trực tiếp vị trí người gặp sự cố. Công trình ``Visionary Vigilance'' xây dựng DiverseFALL10500 với khoảng 10.500 ảnh được chú thích trong nhiều điều kiện chiếu sáng, môi trường, độ tuổi và trang phục; mô hình YOLOv8s được điều chỉnh bằng mô-đun Focus và cơ chế chú ý không gian--kênh CBAM để tăng khả năng biểu diễn \reportcite{ref:falldetection}. Đây là hướng phù hợp khi thiết bị cần quyết định nhanh từ camera RGB.

Ngoài detector theo hộp, Núñez-Marcos, Azkune và Arganda-Carreras biểu diễn chuyển động bằng chồng các ảnh optical flow rồi đưa vào CNN dựa trên VGG-16 để quyết định một đoạn khung hình có chứa cú ngã hay không \reportcite{ref:opticalflow_fall}. Quy trình huấn luyện ba bước bắt đầu từ trọng số ImageNet, học đặc trưng chuyển động trên UCF101 rồi tinh chỉnh cho dữ liệu té ngã, nhằm ứng phó số mẫu ngã ít. Optical flow loại bớt phụ thuộc màu và nền, đồng thời mang thông tin động học mà ảnh tĩnh thiếu; ngược lại, nó tốn bước tính dòng quang học và nhạy với chuyển động camera. Hạn chế sau đặc biệt quan trọng với robot di động, vì chuyển động quay của robot có thể bị hiểu thành chuyển động của người nếu không bù chuyển động nền.

Tuy nhiên, một tư thế nằm trong một ảnh chưa đủ để kết luận đã xảy ra té ngã. Người dùng có thể nằm nghỉ, cúi nhặt đồ, ngồi thấp, tập thể dục hoặc chỉ bị camera quan sát ở góc bất lợi. Detector một khung hình cũng không mô tả được ba pha quan trọng của sự kiện: cơ thể mất thăng bằng, hạ thấp nhanh và duy trì trạng thái bất thường sau đó. Vì vậy, nếu dùng detector ảnh, hệ thống vận hành vẫn cần bộ đệm thời gian, xác nhận liên tiếp và cơ chế hủy cảnh báo khi bằng chứng không bền vững.

\paragraph{Phát hiện dựa trên pose và mô hình chuỗi thời gian.}
Ước lượng tư thế biến mỗi người thành tập điểm khớp $(x,y,c)$, trong đó $c$ là độ tin cậy. Sau khi chuẩn hóa theo hộp người hoặc kích thước cơ thể, có thể suy ra góc thân, tỷ lệ chiều cao--chiều rộng, vận tốc hông/vai, độ cao tương đối và mức thay đổi giữa các khung hình. Luồng xử lý tổng quát là \textit{phát hiện người $\rightarrow$ ước lượng pose $\rightarrow$ gán pose vào đúng track $\rightarrow$ tạo chuỗi đặc trưng $\rightarrow$ LSTM/Transformer $\rightarrow$ làm mượt theo thời gian $\rightarrow$ xác nhận sự kiện}. So với phân lớp ảnh tĩnh, cách này mô hình hóa được diễn biến của cú ngã và làm giảm ảnh hưởng của màu quần áo hoặc nền cảnh.

Nghiên cứu của Juraev và cộng sự đánh giá các phương pháp pose trên dữ liệu hoạt động thực tế của người cao tuổi, đồng thời chỉ ra ba vấn đề nổi bật là thiếu dữ liệu té ngã quy mô lớn, thay đổi góc camera/ánh sáng và che khuất. Nhóm tác giả bổ sung dữ liệu tổng hợp và cho thấy mô hình Transformer nhận dạng hành động tốt hơn LSTM trong thiết lập được khảo sát \reportcite{ref:pose_synthetic}. Kết quả này gợi ý rằng chất lượng keypoint và tính đa dạng dữ liệu quan trọng không kém lựa chọn bộ phân loại chuỗi. Mặc dù Transformer có khả năng mô hình hóa quan hệ dài, LSTM vẫn là lựa chọn hợp lý cho nguyên mẫu di động khi chuỗi ngắn, yêu cầu bộ nhớ thấp và cần chuyển đổi thuận lợi sang định dạng chạy trên thiết bị \reportcite{ref:lstm}.

CAUCAFall là nguồn dữ liệu đáng chú ý vì được thu trong môi trường ít kiểm soát, bao gồm hoạt động thường ngày và các trường hợp ngã với khác biệt về khoảng cách, hướng ngã và độ chiếu sáng; dữ liệu có chú thích phục vụ cả phát hiện người và biểu diễn xương \reportcite{ref:caucafall}. Điều này gần với bài toán triển khai hơn các bộ dữ liệu chỉ quay một phông nền cố định. Dù vậy, các cú ngã trong bộ dữ liệu thường vẫn phải được dàn dựng vì lý do đạo đức và an toàn. Do đó, chia tập ngẫu nhiên theo khung hình có thể làm cùng một người hoặc cùng một cảnh xuất hiện ở cả tập huấn luyện và kiểm thử, tạo kết quả lạc quan. Đánh giá đáng tin cậy nên tách theo người, video hoặc bối cảnh và kiểm thử bổ sung trên camera thực của robot.

\paragraph{Từ dự đoán khung hình đến xác nhận sự kiện.}
Trong chăm sóc tại nhà, đơn vị cần đánh giá là một \emph{sự kiện}, không chỉ một khung hình. Một giải pháp hoàn chỉnh thường giữ cửa sổ dự đoán, yêu cầu đủ số khung hình đồng thuận, kiểm tra thời gian nằm bất động và áp dụng khoảng chờ để không phát nhiều cảnh báo cho cùng một cú ngã. Khi mất keypoint do che khuất, hệ thống không nên tự động diễn giải ``không thấy người'' thành ``an toàn''. Sau xác nhận, pipeline cần chuyển sang trạng thái phản ứng: dừng robot, ghi nhận đúng một sự cố, hỏi người dùng có cần trợ giúp, đặt thời hạn chờ và thông báo người chăm sóc nếu không có phản hồi rõ ràng. Cách nhìn theo sự kiện này làm xuất hiện thêm các chỉ số quan trọng như độ trễ cảnh báo, tỷ lệ báo động giả mỗi giờ, tỷ lệ bỏ sót sự kiện và số cảnh báo trùng, bên cạnh precision, recall và F1 ở mức khung hình.

Việc so sánh kết quả giữa các công trình cần đặc biệt thận trọng vì chúng có thể dùng loại cảm biến, định nghĩa cửa sổ sự kiện, hoạt động thường ngày, đối tượng và cách chia dữ liệu khác nhau. Newaz và Hanada xem thuật toán, bộ dữ liệu và phương pháp validation là ba thành phần phải được báo cáo cùng nhau \reportcite{ref:newaz_fall_review}. Do đó, một giá trị accuracy cao trên các đoạn ngã dàn dựng không đủ chứng minh hệ thống sẽ có ít báo động giả trong nhiều giờ sinh hoạt thật; báo cáo của đồ án ưu tiên trình bày rõ protocol và giới hạn thay vì so sánh trực tiếp các con số không cùng điều kiện.

Đồ án chọn hướng pose--chuỗi thời gian: keypoint được ghép với người đang theo dõi, chuẩn hóa và gom thành cửa sổ 15 khung hình với 69 đặc trưng để đưa vào LSTM. Kết quả mạng không được dùng như một mệnh lệnh đơn lẻ mà đi qua hậu kiểm theo bằng chứng và trạng thái sự cố. Lựa chọn này kế thừa ưu điểm của mô hình pose theo thời gian, đồng thời phù hợp với giới hạn tính toán của smartphone và yêu cầu giảm nhầm lẫn giữa nằm chủ động với một cú ngã thật.

\fi

\subsubsection{Các giải pháp theo dõi người và bám mục tiêu}

Phát hiện người xác định vị trí của đối tượng trong từng khung hình, còn theo dõi tạo sự liên tục giữa các kết quả phát hiện theo thời gian. Trong đề tài, thông tin này được sử dụng với hai mục đích chính: lựa chọn người cần giám sát và xác định sai lệch của người đó so với tâm ảnh để robot điều chỉnh chuyển động. Bài toán được đặt trong môi trường gia đình, nơi phần lớn thời gian chỉ có một người cao tuổi trong vùng quan sát và robot di chuyển ở tốc độ thấp. Vì vậy, yêu cầu cốt lõi không phải là duy trì đồng thời nhiều định danh trong thời gian dài, mà là duy trì ổn định người đã được chọn làm mục tiêu và dừng an toàn khi kết quả quan sát không còn đủ tin cậy.

\paragraph{SORT.}
SORT là phương pháp \textit{tracking-by-detection}, tức sử dụng kết quả của bộ phát hiện làm quan sát đầu vào cho bộ theo dõi. Bộ lọc Kalman dự đoán vị trí hộp ở khung hình kế tiếp; độ chồng lấp IoU biểu diễn mức tương đồng về vị trí giữa hộp dự đoán và hộp vừa phát hiện; thuật toán Hungarian sau đó tìm cách ghép các cặp phù hợp \reportcite{ref:sort}\reportcite{ref:kalman}\reportcite{ref:hungarian}. Phương pháp có cấu trúc gọn, không cần huấn luyện thêm mạng nhận dạng và đáp ứng tốt yêu cầu xử lý thời gian thực trên điện thoại. Hạn chế chính là khả năng giữ ID suy giảm khi mục tiêu bị che khuất lâu, camera chuyển động mạnh hoặc nhiều người đi gần nhau.

Các phương pháp như Deep SORT và ByteTrack cải thiện một số hạn chế trên theo những hướng khác nhau. Deep SORT bổ sung đặc trưng ngoại hình để nhận biết lại đối tượng sau khi bị che khuất, nhờ đó giảm hiện tượng đổi ID trong cảnh có nhiều người \reportcite{ref:deepsort}. ByteTrack tận dụng cả những hộp có độ tin cậy thấp trong bước liên kết, thay vì loại bỏ chúng ngay từ đầu, nên có thể duy trì quỹ đạo tốt hơn khi đối tượng chỉ bị che một phần \reportcite{ref:bytetrack}. Những cải tiến này có ý nghĩa rõ rệt đối với bài toán theo dõi đa đối tượng, nhưng đồng thời làm tăng chi phí tính toán hoặc độ phức tạp của cơ chế quản lý track.

Đối với tình huống một người chuyển nhanh từ tư thế đứng sang nằm, tracker phức tạp cũng không mặc nhiên cho kết quả tốt hơn. Sự thay đổi lớn về chiều cao, tỉ lệ và vị trí của bounding box có thể làm giảm độ tương đồng hình học; trong khi đặc trưng ngoại hình cũng kém ổn định khi góc nhìn và phần cơ thể quan sát được thay đổi đột ngột. Quan trọng hơn, các tracker chỉ giải quyết bài toán liên kết đối tượng, không xác định ý nghĩa của biến đổi tư thế. Do đó, việc nhận biết một cú ngã cần dựa vào keypoint và diễn biến theo thời gian; tracker chỉ hỗ trợ bảo đảm chuỗi tư thế được thu thập từ cùng một người.

\paragraph{Từ theo dõi đến bám người.}
Trong hệ thống của đề tài, người xuất hiện trong vùng quan sát được đánh giá dựa trên kích thước vùng bao, độ tin cậy của phép phát hiện, số keypoint nhìn thấy và vị trí trong khung hình. Người có chất lượng quan sát tốt nhất được chọn làm mục tiêu chính và được duy trì trong các khung hình tiếp theo. Vị trí tâm của bounding box cho biết độ lệch theo phương ngang để robot điều chỉnh hướng quay, còn kích thước vùng bao kết hợp với khoảng cách giữa các keypoint trên cơ thể được dùng để ước lượng mức độ gần--xa. Từ hai đại lượng này, hệ thống xác định thành phần chuyển động tiến--lùi và thành phần điều hướng, sau đó làm mượt lệnh nhằm hạn chế rung khi mục tiêu dao động quanh vùng cân bằng.

Luồng xử lý có thể khái quát thành \textit{thu nhận hình ảnh $\rightarrow$ phát hiện người và keypoint $\rightarrow$ liên kết với mục tiêu chính $\rightarrow$ ước lượng sai lệch hướng và khoảng cách $\rightarrow$ tạo lệnh chuyển động}. Lệnh được ghi nhận trên kênh điều khiển của hệ thống và chuyển xuống thân robot để thực hiện. Cách phân chia này kế thừa tư tưởng của OpenBot, trong đó điện thoại đảm nhiệm cảm nhận và tính toán, còn phần thân robot thực thi lệnh chuyển động \reportcite{ref:openbot}.

Phương pháp theo dõi được xây dựng trên nền tảng SORT nhưng có bổ sung thông tin về khoảng cách tâm, vị trí tương đối của các keypoint và sự thay đổi kích thước bounding box. Histogram màu chỉ được sử dụng khi cần phân biệt giữa nhiều ứng viên hoặc nhận dạng lại mục tiêu sau khi mất dấu, do đó không phải chạy thêm một mạng Re-ID sâu ở mọi khung hình. Cơ chế này giúp giảm nguy cơ đổi ID khi người chuyển từ đứng sang nằm, bởi sự thay đổi mạnh về hình dạng vùng bao không còn được đánh giá chỉ bằng IoU. Nếu ID vẫn bị thay đổi, hệ thống tìm lại mục tiêu trong vùng lân cận của vị trí cuối, yêu cầu kết quả ổn định qua nhiều khung hình rồi mới khôi phục ID chính và tiếp tục chuỗi dữ liệu dùng cho phát hiện té ngã.

Trong chế độ bám thông thường, hệ thống dành một khoảng thời gian ngắn để tái liên kết khi mục tiêu tạm thời biến mất; nếu không tìm lại được trong khoảng hai giây, lệnh dừng được gửi tới robot. Khi đã ghi nhận người đang nằm hoặc đang xử lý một tình huống nghi ngờ té ngã, yêu cầu an toàn được đặt cao hơn: robot dừng ngay khi mất mục tiêu và không tiếp tục tiến về phía trước. Nếu người vẫn còn trong ảnh nhưng vùng quan sát chưa phù hợp, robot chỉ lùi chậm hoặc quay tại chỗ để đưa cơ thể vào gần tâm khung hình, sau đó dừng lại. Như vậy, tracker hỗ trợ duy trì đúng người và cung cấp dữ liệu cho điều khiển, còn kết luận té ngã vẫn do quá trình phân tích tư thế và chuỗi thời gian đảm nhiệm. Bảng~\ref{tab:tracking_related} tóm tắt các lựa chọn được cân nhắc.

\begin{table}[H]
\centering
\caption{So sánh các phương pháp theo dõi người liên quan đến đề tài}
\label{tab:tracking_related}
\small
\begin{tabularx}{\textwidth}{|p{2.2cm}|p{3.0cm}|X|X|}
\hline
\textbf{Phương pháp} & \textbf{Cơ sở liên kết} & \textbf{Ưu điểm và giới hạn} & \textbf{Khả năng áp dụng} \\
\hline
SORT & Kalman, IoU, Hungarian & Gọn, nhanh, không cần mạng Re-ID; dễ đổi ID khi che khuất hoặc giao nhau & Phù hợp làm lõi tracker trên smartphone \\
\hline
Deep SORT & Chuyển động và embedding ngoại hình & Giảm đổi ID; cần suy luận đặc trưng và tài nguyên lớn hơn & Hữu ích nếu thiết bị đủ mạnh và nhiều người xuất hiện \\
\hline
ByteTrack & Ghép hai lượt với hộp điểm cao và thấp & Khôi phục quan sát bị che khuất; nhạy với ngưỡng và chất lượng detector & Gợi ý giữ bằng chứng yếu có kiểm soát \\
\hline
Hướng của đề tài & SORT mở rộng với hình học, keypoint và histogram màu & Duy trì mục tiêu qua thay đổi tư thế mà không cần mạng Re-ID sâu; khả năng phục hồi vẫn phụ thuộc chất lượng ảnh & Phù hợp với bối cảnh một mục tiêu và xử lý thời gian thực trên điện thoại \\
\hline
\end{tabularx}
\end{table}

\subsubsection{Chatbot AI và hội thoại hỗ trợ người cao tuổi}

Trong đồ án, chatbot được định hướng như một người bạn trò chuyện bằng giọng nói trên robot. Người dùng có thể gọi robot, đặt câu hỏi, kể chuyện hoặc trao đổi về những chủ đề hằng ngày mà không phải thao tác với bàn phím và màn hình nhỏ. Mục tiêu chính là tạo một kênh giao tiếp tự nhiên, phản hồi ngắn gọn và dễ sử dụng; chatbot không được xem là công cụ chẩn đoán, trị liệu tâm lý hay tư vấn y khoa.

\paragraph{Kiến trúc của một chatbot giọng nói.}
Một lượt hội thoại hoàn chỉnh thường đi qua chuỗi \textit{kích hoạt $\rightarrow$ thu âm $\rightarrow$ nhận dạng tiếng nói (STT) $\rightarrow$ mô hình ngôn ngữ lớn (LLM) $\rightarrow$ tổng hợp tiếng nói (TTS) $\rightarrow$ phát âm thanh}. Bên cạnh ba khối xử lý chính, ứng dụng còn cần phát hiện khoảng lặng, quản lý lịch sử hội thoại, ngắt câu trả lời đang phát, xử lý mất mạng và biểu diễn trạng thái để người dùng biết robot đang nghe, suy nghĩ hay nói. Chất lượng trải nghiệm vì thế phụ thuộc vào độ trễ đầu cuối và cách phối hợp các thành phần, không chỉ vào chất lượng văn bản do LLM sinh ra.

Xiaozhi-ESP32 là một hiện thực tham khảo tiêu biểu cho cách tổ chức này trên thiết bị có tài nguyên hạn chế. Nền tảng thực hiện đánh thức bằng từ khóa ngay trên thiết bị, sau đó truyền âm thanh nén Opus qua WebSocket hoặc MQTT/UDP để sử dụng chuỗi ASR--LLM--TTS; giao thức còn quy định các trạng thái bắt đầu nghe, bắt đầu nói, kết thúc nói và hủy lượt hội thoại \reportcite{ref:xiaozhi}. Điểm phù hợp với đồ án không nằm ở phần cứng ESP32 mà ở nguyên tắc kiến trúc: thiết bị đầu cuối chỉ duy trì vòng đời phiên, thu/phát âm thanh và giao diện trạng thái, còn tác vụ ngôn ngữ nặng có thể được chuyển tới dịch vụ bên ngoài. Cách phân vai này cho phép một thiết bị nhỏ tạo cảm giác hội thoại liên tục mà không phải chứa toàn bộ mô hình AI.

\paragraph{Các lựa chọn cho khối STT và TTS.}
Các giải pháp tiếng nói có thể được chia thành hai nhóm theo vị trí xử lý. Nhóm thứ nhất thực hiện STT và TTS ngay trên thiết bị. Android cung cấp khả năng nhận dạng cục bộ trên những thiết bị có dịch vụ tương thích, đồng thời có thể tổng hợp giọng nói từ các gói giọng đã được cài đặt \reportcite{ref:android_speech}\reportcite{ref:android_tts}. Ưu điểm của lựa chọn này là ứng dụng không phải đóng gói thêm mô hình, độ trễ thấp và có thể hoạt động khi mất mạng. Tuy nhiên, mức hỗ trợ tiếng Việt và chất lượng nhận dạng phụ thuộc vào phiên bản hệ điều hành, nhà sản xuất và dữ liệu ngôn ngữ có sẵn; dịch vụ mặc định trên một số thiết bị vẫn có thể cần kết nối mạng. Vì vậy, STT/TTS của Android nên được xem là giải pháp có khả năng xử lý cục bộ, thay vì một cơ chế ngoại tuyến đồng nhất trên mọi điện thoại.

Khi cần chủ động hơn về khả năng ngoại tuyến, mô hình tiếng nói có thể được tích hợp cùng ứng dụng. Vosk hỗ trợ nhận dạng luồng và cung cấp các mô hình nhỏ dành cho thiết bị di động, trong đó có tiếng Việt \reportcite{ref:vosk_android}. Sherpa-onnx có phạm vi rộng hơn, bao gồm nhận dạng tiếng nói, phát hiện hoạt động giọng nói, nhận biết từ khóa và tổng hợp tiếng nói trên Android \reportcite{ref:sherpa_onnx}. Hai nền tảng này giữ dữ liệu âm thanh trên thiết bị và không phát sinh chi phí theo lượt sử dụng, nhưng làm tăng dung lượng cài đặt, mức sử dụng CPU/RAM và công việc lựa chọn mô hình. Chất lượng thực tế cũng cần được đánh giá trong điều kiện có tiếng động cơ và tạp âm sinh hoạt trong nhà.

Nhóm thứ hai sử dụng dịch vụ tiếng nói trực tuyến. Âm thanh được gửi theo từng đoạn ngắn lên máy chủ; dịch vụ STT trả dần kết quả tạm thời và kết quả hoàn chỉnh, trong khi TTS có thể truyền các khối âm thanh về để phát trước khi toàn bộ câu được tổng hợp xong. Azure Speech là một ví dụ cung cấp SDK cho luồng âm thanh thời gian thực và hỗ trợ cả nhận dạng lẫn giọng đọc tiếng Việt \reportcite{ref:azure_speech_sdk}\reportcite{ref:azure_speech_languages}. Cách tiếp cận này có ưu thế về chất lượng mô hình, khả năng cập nhật và độ nhất quán giữa các thiết bị. Tuy nhiên, quá trình xử lý phụ thuộc vào đường truyền và thường phát sinh chi phí theo thời lượng âm thanh hoặc lượng văn bản sau khi vượt phạm vi dùng thử. Trong điều kiện thời gian và kinh phí của đồ án, hướng trực tuyến cho STT/TTS được xem như một phương án tham khảo, còn giải pháp sẵn có trên thiết bị thuận lợi hơn cho việc xây dựng nguyên mẫu.

\paragraph{Lựa chọn giữa LLM qua API và LLM chạy trên thiết bị.}
LLM qua API cho phép điện thoại gửi nội dung đã nhận dạng cùng một phần ngữ cảnh hội thoại và nhận lại phản hồi theo từng đoạn. Khả năng streaming giúp ứng dụng bắt đầu xử lý câu trả lời sớm, nhờ đó giảm thời gian chờ cảm nhận của người dùng. Phương án này giữ ứng dụng gọn, tận dụng được mô hình có năng lực tốt hơn phần cứng di động và cho phép thay đổi mô hình tương đối linh hoạt. Hạn chế chủ yếu là sự phụ thuộc vào mạng và chi phí có thể phát sinh nếu nhu cầu sử dụng trong quá trình thực hiện đồ án vượt quá giới hạn của gói miễn phí.

Groq cung cấp API hội thoại và phản hồi streaming, phù hợp để khảo sát chatbot có độ trễ thấp \reportcite{ref:groq_api}. Gói miễn phí giới hạn số yêu cầu và token nhưng vẫn đáp ứng nhu cầu xây dựng, điều chỉnh và đánh giá nguyên mẫu \reportcite{ref:groq_limits}. Gemini Developer API cũng cung cấp một số mô hình có thể sử dụng miễn phí trong phạm vi hạn mức nhất định \reportcite{ref:gemini_pricing}.

Ở hướng ngược lại, LLM cục bộ giữ nội dung hội thoại trên điện thoại và không phụ thuộc mạng sau khi mô hình đã được tải xuống. Google AI Edge Gallery và MLC LLM cho thấy các mô hình đã lượng tử hóa có thể được triển khai trên Android \reportcite{ref:ai_edge_gallery}\reportcite{ref:mlc_android}. Dù vậy, mô hình vẫn chiếm đáng kể dung lượng lưu trữ và bộ nhớ, đồng thời cạnh tranh tài nguyên với pipeline thị giác của robot. Việc giảm kích thước mô hình còn có thể ảnh hưởng đến khả năng hiểu tiếng Việt và chất lượng hội thoại. Vì vậy, LLM cục bộ phù hợp hơn với thiết bị có cấu hình đủ mạnh hoặc yêu cầu ngoại tuyến và riêng tư cao; LLM qua API phù hợp hơn với nguyên mẫu cần hội thoại tự nhiên nhưng có tài nguyên tính toán giới hạn.

\begin{table}[H]
\centering
\caption{So sánh các phương án tạo phản hồi hội thoại}
\label{tab:android_chatbot_options}
\small
\begin{tabularx}{\textwidth}{|p{2.6cm}|X|X|X|}
\hline
\textbf{Phương án} & \textbf{Ưu điểm} & \textbf{Hạn chế} & \textbf{Mức phù hợp} \\
\hline
Luật/kho câu trả lời & Rất nhẹ, ngoại tuyến, kết quả xác định & Hội thoại cứng và khó bao phủ cách diễn đạt tự nhiên & Tương tác ngắn và nội dung chuẩn bị trước \\
\hline
LLM qua API & Ứng dụng nhẹ, phản hồi tốt, hỗ trợ streaming, dễ đổi mô hình & Cần mạng; chi phí tăng khi vượt giới hạn miễn phí & Phù hợp với nguyên mẫu hội thoại hằng ngày \\
\hline
LLM cục bộ & Ngoại tuyến, dữ liệu không rời thiết bị, không tính phí theo lượt & Mô hình lớn, dùng nhiều RAM và năng lượng; phụ thuộc phần cứng & Hướng phát triển khi thiết bị đáp ứng tài nguyên \\
\hline
\end{tabularx}
\end{table}

\paragraph{Tổ chức hội thoại phù hợp với người cao tuổi.}
Một chatbot trò chuyện hằng ngày không cần câu trả lời dài hoặc lập luận phức tạp. Giá trị sử dụng đến từ việc kích hoạt dễ nhớ, nghe được nhiều cách diễn đạt, trả lời đúng ngôn ngữ, duy trì một lượng ngữ cảnh vừa đủ và phản hồi sớm. Giao diện nên thể hiện rõ các trạng thái đang nghe, đang xử lý và đang nói; cho phép lặp lại khi STT nhận sai; tránh đọc ký hiệu, đường dẫn hoặc đoạn văn quá dài. Hồ sơ như tên gọi, cách xưng hô và sở thích có thể giúp cuộc trò chuyện tự nhiên hơn, nhưng chỉ nên đưa vào prompt những trường cần thiết cho lượt hội thoại hiện tại.

Ba phương án trong bảng không loại trừ lẫn nhau. LLM thích hợp với hội thoại mở vì có thể tạo phản hồi linh hoạt, trong khi các câu thông báo và tương tác ngắn có thể sử dụng nội dung được chuẩn bị trước. Việc phối hợp hai cách tiếp cận giúp giới hạn vai trò của mô hình ngôn ngữ ở những phần cần khả năng trò chuyện, thay vì dùng nó cho mọi phản hồi của ứng dụng. Cách phân chia này phù hợp với nhận định từ các nghiên cứu về conversational agent rằng khả năng tiếp nhận ngôn ngữ tự nhiên không đồng nghĩa với việc toàn bộ hoạt động của hệ thống phải do mô hình sinh điều khiển \reportcite{ref:healthcare_chatbot_review}.

\paragraph{Hướng lựa chọn cho đề tài.}
Với mục tiêu tạo một người bạn trò chuyện bằng tiếng Việt trên điện thoại Android, đề tài lựa chọn các dịch vụ nhận dạng và tổng hợp tiếng nói sẵn có trên thiết bị kết hợp với LLM qua API. Phương án này không yêu cầu đóng gói thêm mô hình STT/TTS hoặc LLM dung lượng lớn, phù hợp với giới hạn tài nguyên của điện thoại đồng thời đáp ứng nhu cầu thử nghiệm hội thoại hằng ngày.

Do LLM được cung cấp qua API, chức năng trò chuyện vẫn phụ thuộc Internet ngay cả khi STT và TTS có thể chạy cục bộ. Việc tách riêng ba khối STT, LLM và TTS tạo điều kiện khảo sát, đánh giá và thay thế từng thành phần bằng giải pháp trực tuyến hoặc cục bộ khi yêu cầu và khả năng phần cứng thay đổi.

\subsubsection{Các nền tảng robot tích hợp và bài học áp dụng}

SocialRobot là một ví dụ về robot di động hỗ trợ chăm sóc theo hướng lấy người dùng làm trung tâm. Nền tảng kết hợp khả năng di chuyển và cảm nhận tại chỗ với hệ thống quản lý hồ sơ, dịch vụ cá nhân hóa và nhóm chăm sóc từ xa \reportcite{ref:socialrobot}. Điểm đáng tham khảo là người thân và người chăm sóc được xem như những tác nhân tham gia vào quá trình hỗ trợ, thay vì chỉ nhận thông báo khi sự cố đã xảy ra.

OpenBot tiếp cận bài toán từ hướng đơn giản hóa nền tảng robot: điện thoại thông minh đảm nhiệm cảm nhận, tính toán, giao diện và kết nối, còn phần thân xe tập trung thực hiện lệnh chuyển động \reportcite{ref:openbot}\reportcite{ref:openbot_repo}. Cách tổ chức này làm giảm số phần cứng chuyên dụng và tạo điều kiện thuận lợi cho học tập, thử nghiệm. Đề tài kế thừa tư tưởng sử dụng điện thoại làm trung tâm xử lý, đồng thời bổ sung ứng dụng dành cho người chăm sóc và các chức năng phù hợp với bối cảnh hỗ trợ người cao tuổi tại nhà.

\subsubsection{Khoảng trống nghiên cứu và định vị của đề tài}

Các công trình được khảo sát thường tập trung giải quyết tốt một bài toán riêng. Nghiên cứu phát hiện té ngã chủ yếu đánh giá mô hình trên bộ dữ liệu; nghiên cứu theo dõi chú trọng duy trì danh tính; chatbot tập trung vào chất lượng tương tác; còn các nền tảng robot giá thấp quan tâm đến khả năng điều khiển và mở rộng. Những nghiên cứu thị giác gần đây đã cải thiện đáng kể khả năng xử lý nhiều người, che khuất và triển khai trên thiết bị biên \reportcite{ref:hfdmia_pose}\reportcite{ref:yosap_lstm}\reportcite{ref:fallnet}, nhưng việc đưa kết quả nhận biết vào một quy trình chăm sóc hoàn chỉnh vẫn cần được xem xét thêm.

Khoảng trống mà đề tài hướng tới là sự phối hợp từ quan sát, nhận biết tình huống đến tương tác và thông báo cho người chăm sóc. Khi các chức năng cùng hoạt động trên robot di động, kết quả của bước trước ảnh hưởng trực tiếp đến bước sau; vì vậy hệ thống không chỉ cần mô hình nhận biết tốt mà còn phải duy trì đúng người mục tiêu, xử lý trường hợp thiếu dữ liệu và chuyển về trạng thái an toàn khi không đủ cơ sở ra quyết định.

\subsection{Đóng góp của đề tài}

Đề tài hướng đến đóng góp ở khả năng tổ chức nhiều chức năng hỗ trợ thành một hệ thống thống nhất, phù hợp với bối cảnh chăm sóc người cao tuổi tại nhà. Những kết quả chính đạt được gồm:

\begin{itemize}
    \item Xây dựng nguyên mẫu robot trợ lý có khả năng hỗ trợ giao tiếp, nhắc nhở, giám sát và kết nối người cao tuổi với người chăm sóc. Sự phối hợp giữa các chức năng này góp phần hình thành một hệ sinh thái chăm sóc liên tục, trong đó nhu cầu sinh hoạt hằng ngày và các tình huống cần trợ giúp được tiếp nhận trên cùng một hệ thống.
    \item Đề xuất cách phân chia chức năng và luồng phối hợp giữa người cao tuổi, robot và người chăm sóc. Kết quả này cho thấy khả năng kết hợp trợ lý ảo, giám sát an toàn và robot di động trong một giải pháp hỗ trợ tại nhà, đồng thời tạo cơ sở để bổ sung các dịch vụ chăm sóc khác trong tương lai.
    \item Xây dựng một nền tảng thử nghiệm phục vụ việc nghiên cứu các bài toán liên quan đến tương tác người--robot, hỗ trợ sinh hoạt, nhận biết tình huống bất thường và phối hợp chăm sóc từ xa. Nguyên mẫu có thể được tiếp tục sử dụng để khảo sát phương pháp mới, so sánh giải pháp hoặc đánh giá trải nghiệm của người dùng trong những nghiên cứu tiếp theo.
    \item Góp phần làm rõ tiềm năng ứng dụng của trí tuệ nhân tạo và robot trong hỗ trợ người cao tuổi, qua đó cung cấp một tài liệu tham khảo cho sinh viên và các nhóm nghiên cứu khi phát triển những ứng dụng có mục tiêu tương tự. Kết quả của đề tài cũng đặt nền tảng cho việc hoàn thiện giải pháp theo hướng dễ tiếp cận, tin cậy và phù hợp hơn với điều kiện chăm sóc thực tế.
\end{itemize}

\subsection{Cấu trúc báo cáo}

Báo cáo được tổ chức thành sáu chương, sắp xếp theo trình tự từ xác định bài toán, xây dựng nền tảng lý thuyết đến thiết kế, hiện thực và đánh giá hệ thống. Nội dung cụ thể của từng chương như sau:

\begin{itemize}
    \item \textbf{Chương 1 -- Giới thiệu:} trình bày bối cảnh già hóa dân số và nhu cầu hỗ trợ người cao tuổi tại nhà; từ đó xác định mục tiêu, đối tượng, phạm vi, phương pháp nghiên cứu, các công trình liên quan và đóng góp của đề tài.
    \item \textbf{Chương 2 -- Cơ sở lý thuyết:} hệ thống hóa những kiến thức làm nền tảng cho việc xây dựng robot trợ lý, bao gồm nền tảng robot sử dụng điện thoại thông minh, thị giác máy tính, mô hình nhận dạng tư thế và chuỗi thời gian, xử lý hội thoại, suy luận trên thiết bị, truyền thông và điều khiển phần thân robot.
    \item \textbf{Chương 3 -- Phân tích và thiết kế hệ thống:} phân tích yêu cầu chức năng và phi chức năng, xây dựng kiến trúc tổng quan, mô tả các luồng hoạt động chính, thiết kế các khối chức năng, tổ chức dữ liệu và những yêu cầu về an toàn, bảo mật.
    \item \textbf{Chương 4 -- Hiện thực hệ thống:} trình bày quá trình xây dựng các thành phần của nguyên mẫu, gồm ứng dụng trên robot, ứng dụng dành cho người chăm sóc, các chức năng trí tuệ nhân tạo, quá trình chuẩn bị mô hình và chương trình điều khiển phần thân robot.
    \item \textbf{Chương 5 -- Kết quả và thảo luận:} trình bày phương pháp và tiêu chí đánh giá, kết quả kiểm thử chức năng, kết quả đánh giá mô hình trí tuệ nhân tạo và hiệu năng hệ thống; đồng thời phân tích những kết quả đạt được và các hạn chế còn tồn tại.
    \item \textbf{Chương 6 -- Kết luận:} tổng kết quá trình thực hiện, đối chiếu kết quả với mục tiêu đề tài và đề xuất các hướng nghiên cứu, hoàn thiện hệ thống trong tương lai.
\end{itemize}

Phần phụ lục bổ sung các đặc tả trường hợp sử dụng, cấu trúc dữ liệu, danh mục thành phần, thông tin mô hình và biểu mẫu kiểm thử nhằm hỗ trợ việc tra cứu và kiểm chứng nội dung được trình bày trong báo cáo.

\subsection{Kết luận chương}

Chương 1 đã làm rõ bối cảnh và tính cấp thiết của việc nghiên cứu giải pháp hỗ trợ chăm sóc người cao tuổi tại nhà. Quá trình già hóa dân số, nguy cơ té ngã, sự suy giảm trí nhớ, nhu cầu duy trì lịch sinh hoạt và khoảng cách giữa người cao tuổi với người chăm sóc đặt ra yêu cầu về một phương thức hỗ trợ chủ động, thuận tiện và có khả năng phản hồi kịp thời. Trong bối cảnh đó, robot trợ lý được xem là một phương tiện bổ trợ cho hoạt động chăm sóc, giúp duy trì tương tác hằng ngày, hỗ trợ giám sát và tạo thêm kênh kết nối với người thân, thay vì thay thế vai trò của con người hoặc các thiết bị y tế chuyên dụng.

Từ bài toán đặt ra, chương đã xác định mục tiêu xây dựng một nguyên mẫu tích hợp các chức năng giao tiếp, nhắc nhở, theo dõi, nhận biết tình huống nghi té ngã và hỗ trợ liên lạc từ xa. Đối tượng, phạm vi chức năng và giới hạn nghiên cứu cũng được xác định nhằm bảo đảm nội dung thực hiện phù hợp với điều kiện của đồ án. Bên cạnh đó, phương pháp nghiên cứu được tổ chức từ khảo cứu tài liệu, phân tích và thiết kế đến xây dựng nguyên mẫu, thực nghiệm và đánh giá. Việc khảo sát các công trình liên quan cho thấy mỗi hướng tiếp cận đều có những ưu điểm và giới hạn riêng; đồng thời cung cấp cơ sở để lựa chọn giải pháp phù hợp với môi trường trong nhà và khả năng triển khai trên robot di động.

Trên cơ sở đó, đóng góp của đề tài được định hướng vào việc kết nối các chức năng hỗ trợ thành một hệ thống thống nhất dành cho người cao tuổi và người chăm sóc, đồng thời hình thành một nền tảng có thể tiếp tục mở rộng trong các nghiên cứu sau. Những vấn đề đã xác lập trong chương này là cơ sở để Chương 2 trình bày các kiến thức lý thuyết cần thiết, trước khi tiến hành phân tích và thiết kế hệ thống ở Chương 3.
